% !TeX root = ./rpz.tex
\Conclusion % заключение к отчёту

В ходе выполнения выпускной квалификационной работы разработано веб-приложение для автоматизированного планирования маршрутов с учётом временных окон.

В результате работы выполнены следующие задачи:

\begin{compactlist}
  \item Проведён анализ предметной области: выявлены ограничения существующих инструментов управления задачами и навигационных сервисов, обоснована актуальность разработки.
  \item Исследованы алгоритмические подходы к решению задачи маршрутизации с временными окнами; выбрана и реализована эвристика ближайшего соседа (NNH-TW) с временной сложностью O(n²).
  \item Разработана архитектура программной системы на основе принципов чистой архитектуры с разделением на слои: доменная логика, варианты использования, HTTP-доставка и репозитории баз данных.
  \item Реализована серверная часть на языке Go с использованием фреймворка Gin, обеспечивающая REST API для управления задачами, маршрутами и аутентификацией пользователей.
  \item Разработан пользовательский интерфейс на React с отображением задач в виде карточек, интерактивной картой Яндекс Карт и возможностью ручного упорядочивания методом перетаскивания.
  \item Реализована интеграция с Яндекс Картами JavaScript API 2.1 для геокодирования адресов и визуализации маршрута, а также с OSRM для формирования матрицы расстояний.
  \item Проведено тестирование разработанного решения: реализованы модульные тесты доменной логики и HTTP-обработчиков.
  \item Выполнена оценка производительности и масштабируемости системы: подтверждено соответствие нефункциональным требованиям по времени ответа, надёжности и устойчивости к ограничениям внешних API.
\end{compactlist}

Поставленная цель достигнута: разработанное веб-приложение устраняет выявленный в ходе анализа разрыв между инструментами управления задачами и навигационными сервисами, предоставляя пользователю единый инструмент для автоматизированного планирования маршрутов с учётом географического положения и временных ограничений.

Разработанная система ориентирована на частных пользователей, выездных специалистов, курьеров и субъектов малого бизнеса, что подтверждается анализом целевой аудитории в первой главе. Применение Docker Compose обеспечивает воспроизводимость развёртывания и независимость от конфигурации целевого сервера.

Перспективами развития системы являются: интеграция с реальными картографическими API для построения реального времени движения, добавление поддержки нескольких пользователей в рамках одного маршрута, разработка мобильного приложения для работы в оффлайн-режиме.
